\clearpage\section{Конструкторская часть}
\subsection{Последовательный алгоритм}
В единственном потоке создаётся нулевая гистограмма, подсчитываются частоты яркостей, определяется порог по формуле~\eqref{eq:criterion} и формируется маска по формуле~\eqref{eq:mask}. На рисунке~\ref{fig:seq} приведена схема для допустимых входных данных; процедура выбора порога общая для обеих версий. Обозначение $[P]$ равно единице, если условие $P$ истинно, иначе нулю.
\diagram{sequential}{Схема алгоритма последовательного потока}{fig:seq}

\clearpage\subsection{Распределение работы}
Число рабочих потоков обозначается $q$. Диспетчер не входит в это число и не обрабатывает собственный участок изображения. Рабочему $j$, $0\le j<q$, назначается полуинтервал $[b_j,e_j)$:
\begin{equation}\label{eq:partition}
\begin{aligned}
d&=\lfloor N/q\rfloor,\quad r=N\bmod q,\\
b_j&=j\cdot d+\min(j,r),\\
e_j&=b_j+d+[j<r].
\end{aligned}
\end{equation}
Размеры участков различаются не более чем на один пиксель. Первый участок начинается с нуля, последний заканчивается на $N$, а конец каждого участка совпадает с началом следующего. Поэтому каждый пиксель обрабатывается ровно одним рабочим. При $q>N$ часть участков пуста, но фактически создаются все заданные $q$ потоков.

Для связи потоков используются мьютекс, то есть средство взаимоисключающего доступа к состоянию синхронизации, и две условные переменные для ожидания изменений этого состояния. Счётчик $c$ хранит число построенных гистограмм, флаг $p$ разрешает бинаризацию, флаг $a$ обозначает отмену запуска. Начальные значения равны нулю. Все обращения к этим переменным и публикация общего порога защищены одним мьютексом.

\subsection{Алгоритм потока-диспетчера}
После выделения памяти под локальные гистограммы и описатели создаются рабочие потоки. Диспетчер ожидает готовности всех гистограмм, объединяет их и вычисляет порог. Затем устанавливается флаг разрешения второй фазы, пробуждаются все ожидающие рабочие и освобождается мьютекс. Завершение каждого созданного потока ожидается до освобождения его данных. При неудачном запуске очередного рабочего устанавливается флаг отмены; уже созданные рабочие пробуждаются и завершаются без бинаризации.

На рисунке~\ref{fig:dispatcher} приведена схема диспетчера. Проверка аргументов и возврат при невозможности выделить или инициализировать служебные данные выполняются перед показанным запуском рабочих.
\diagram{dispatcher}{Схема алгоритма потока-диспетчера}{fig:dispatcher}

\clearpage\subsection{Алгоритм рабочего потока}
Используется один вид рабочего потока. Каждый рабочий получает свой участок изображения и маски, а также обнулённую диспетчером гистограмму. После подсчёта частот увеличивается $c$, уведомляется диспетчер и ожидается порог либо отмена. Затем независимо формируется участок маски. На рисунке~\ref{fig:worker} приведена схема рабочего потока.
\diagram{worker}{Схема алгоритма рабочего потока}{fig:worker}

\clearpage\subsection{Взаимодействие и корректность}
На рисунке~\ref{fig:interaction} представлен порядок взаимодействия потоков. Уведомления о готовности поступают после завершения соответствующей локальной гистограммы; второй этап разрешается только после получения всех уведомлений и вычисления общего порога.
\diagram{interaction}{Порядок взаимодействия диспетчера и рабочих потоков}{fig:interaction}

Ожидание условной переменной временно освобождает мьютекс; после пробуждения он вновь захватывается~\cite{posix}. Предикат ожидания проверяется в цикле, поэтому ложное пробуждение не разрешает преждевременную бинаризацию. После копирования порога мьютекс освобождается.

Локальные гистограммы принадлежат разным рабочим, поэтому при подсчёте частот конфликтов записи нет. Для каждого уровня яркости сумма локальных счётчиков совпадает со счётчиком последовательной версии. Суммирование целочисленное, порядок объединения результата не изменяет. В обеих версиях применяется одна процедура определения порога, включая правило выбора при равенстве критерия.

Видимость законченных гистограмм и вычисленного порога обеспечивается освобождением и захватом мьютекса. Для уведомлений о готовности и о продолжении используются разные условные переменные. Маска разбита на непересекающиеся участки; общий вход во время вычисления не изменяется. После ожидания завершения всех рабочих выходная маска полностью сформирована и побайтно совпадает с последовательным результатом.

\subsection{Оценка затрат ресурсов}
Пусть $L=256$ --- число уровней яркости. Последовательная трудоёмкость составляет $O(N+L)$, дополнительная память без учёта входа и маски --- $O(L)$. Суммарная работа параллельной версии равна $O(N+q\cdot L)$, дополнительная память --- $O(q\cdot L+q)$. Для маски в обеих версиях требуется ещё $N$ байт.

При одновременном исполнении всех рабочих и равномерной загрузке оценка вычислительной части критического пути равна
\begin{equation}
O\bigl(\lceil N/q\rceil+q\cdot L+L\bigr).
\end{equation}
К реальному времени добавляются расходы на выделение памяти, запуск, синхронизацию и завершение потоков. При $q$ больше числа доступных логических ядер все рабочие одновременно не исполняются, поэтому член $N/q$ не описывает фактическое сокращение времени. Объединение гистограмм и поиск порога остаются последовательными этапами.

\subsection{Требования к программе и тестирование}
Программа читает и записывает двоичные полутоновые изображения PGM с сигнатурой P5 и максимальным значением 255. Нулевые размеры, переполнение размера, усечённый растр, лишние байты и неподдерживаемый формат отклоняются. Число рабочих потоков задаётся явно в диапазоне от 1 до 4096. В библиотечных функциях допускается использование одного буфера для входа и результата; частичное перекрытие буферов не допускается.

Для проверки выбраны изображения с заранее определяемыми порогом и маской: однотонные, бинарные, с двумя различными яркостями и со всеми уровнями серого. В обоих режимах сравнивались порог и значения пикселей полученной маски. Примеры проверок приведены в таблице~\ref{tab:tests}; векторы рассматриваются как изображения высотой один пиксель.
\begin{table}[H]\caption{Примеры функциональных проверок}\label{tab:tests}
\centering\begin{tabularx}{\textwidth}{|X|c|X|}\hline
Вход & Порог & Маска / результат\\\hhline{|=:=:=|}
$(0)$ & 0 & $(0)$\\\hline
$(255)$ & 255 & $(0)$\\\hline
15 пикселей яркости 77 & 77 & все значения 0\\\hline
$(0,255)$ & 0 & $(0,255)$\\\hline
$(10,10,200,200)$ & 10 & $(0,0,255,255)$\\\hline
$(0,2,4)$ & 0 & $(0,255,255)$\\\hline
$(0,1,\ldots,255)$ & 127 & 128 нулей, 128 значений 255\\\hline
Нулевой размер, усечённый PGM & --- & ошибка ввода\\\hline
$q=0$, $q=-1$, $q=4097$ & --- & ошибка аргумента\\\hline
\end{tabularx}\end{table}

На однотонном изображении проверялось получение нулевой маски, на бинарном --- сохранение разделения на чёрные и белые пиксели. Для входа $(0,2,4)$ проверялся выбор меньшего порога при равенстве критерия.

Последовательная и параллельная версии дали одинаковые результаты. Дополнительно проверялись случаи, когда размер изображения не делится на число потоков и когда потоков больше, чем пикселей. При неверном числе потоков, нулевом размере и усечённом изображении программа выдавала сообщение об ошибке.
